\section{Related Work}
\label{sec:related}

\paragraph{Blockchain Governance and Self-Amendment.}
The problem of protocol upgrades is central to blockchain
systems. Tezos~\cite{goodman2014} introduced self-amending ledgers
where stakeholders vote on protocol changes through an on-chain
governance mechanism. However, Tezos and similar systems rely on
\emph{token-weighted voting}: upgrade authority is proportional to
economic stake, which is a form of plutocracy rather than
democracy~\cite{buterin2021}. This maps poorly to the Separation of
Powers required for algorithmic institutions governing individual
rights: in such contexts, the operator (with the largest economic
stake) would have the most influence over the rules constraining
the operator---a direct inversion of the accountability goal.
Constitutional Code replaces stake-weighted voting with
\emph{branch-weighted ratification}: authority is distributed across
the three constitutional branches, regardless of economic stake.

\paragraph{Regulatory Capture.}
Stigler~\cite{stigler1971} established that regulatory agencies tend
over time to be captured by the industries they regulate, because
the regulated industries have concentrated economic incentives to
influence regulation while the public has diffuse incentives to
resist it. Technical proposals such as multi-signature wallets
mitigate single-point-of-failure risks but are insufficient if all
keyholders belong to the same equilibrium of interests---as would
be the case in a regulator-operator coalition. The Tripartite
Threshold proposed in \S\ref{subsec:threshold} forces a
\emph{cross-branch} coalition for constitutional amendments,
making capture significantly more expensive: an adversary must
simultaneously compromise principals whose interests are
structurally opposed.

\paragraph{Sunset Legislation.}
Sunset clauses in legal systems ensure that certain laws expire
unless renewed by positive legislative action~\cite{kerwin2018}.
They have been applied in administrative law, fiscal policy, and
intelligence authorisation. Their purpose---forcing periodic
democratic recommitment---is well-established; their failure mode
is equally well-documented: legislative inertia leads to automatic
renewal without substantive review~\cite{graetz2010}. This paper
operationalises sunset clauses as \emph{linear resource types},
enforcing expiration at the type-system level. Unlike statutory
sunset clauses, which rely on legislative memory to not re-authorise,
our mechanism relies on the impossibility of constructing a
\texttt{Verdict} without a live \texttt{Mandate} token: the system
cannot forget to expire.

\paragraph{Constitutional Amendment Theory.}
Si\`{e}yes's distinction between \emph{pouvoir constituant} and
\emph{pouvoir constitu\'{e}}~\cite{sieyes1789} is the foundational
concept of constitutional theory: the power that creates the
constitution is categorically different from the powers exercised
under it. This paper is, to our knowledge, the first formal
cryptographic implementation of this distinction. The Stone Clauses
(Definition~\ref{def:amendment-classes}) are the type-theoretic
realisation of \emph{pouvoir constituant}: they cannot be altered
by any power that the constitution itself created.
